\documentclass[a4paper,12pt]{article}
\usepackage[margin=2.5cm]{geometry}
\usepackage{amsmath,amssymb,amsfonts}
\usepackage{tcolorbox}
\usepackage{booktabs}
\usepackage{enumerate}
\usepackage{setspace}
\usepackage{xcolor}

\usepackage{xepersian}
\settextfont{XB Niloofar}
\setdigitfont{XB Niloofar}

% تنظیم فاصله بین خطوط به ۱.۲۵ برابر برای خوانایی بهتر
\setstretch{1.25}

\begin{document}

\begin{tcolorbox}[colframe=black, colback=white, sharp corners, boxrule=1pt, arc=3pt, left=15pt, right=15pt, top=10pt, bottom=10pt]
\noindent
\begin{tabular}{@{}p{0.55\textwidth} @{}p{0.45\textwidth}}
\textbf{درس:} داده‌کاوی & \textbf{دانشکده:} علوم ریاضی -- گروه علوم کامپیوتر \\
\textbf{تمرین:} سری سوم: سیستم‌های توصیه‌گر & \textbf{استاد:} دکتر سهیلا اشکذری طوسی \\
\textbf{ترم:} بهار ۱۴۰۵ & \textbf{دانشگاه:} فردوسی مشهد \\
\end{tabular}
\end{tcolorbox}

\vspace{0.5cm}

\begin{enumerate}

\item \textbf{محاسبه دستی فیلترینگ مشارکتی کاربر‌محور (\lr{User-Based CF})}
\\
ماتریس امتیازات زیر را برای ۴ کاربر و ۴ فیلم در نظر بگیرید. علامت سؤال ($?$) نشان‌دهنده امتیازات مفقود است:

\begin{table}[h!]
\centering
\begin{tabular}{ccccc}
\toprule
\textbf{کاربر / فیلم} & \textbf{فیلم ۱} & \textbf{فیلم ۲} & \textbf{فیلم ۳} & \textbf{فیلم ۴} \\
\midrule
\textbf{کاربر ۱ ($U_1$)} & $5$ & $3$ & $?$ & $1$ \\
\textbf{کاربر ۲ ($U_2$)} & $4$ & $?$ & $4$ & $2$ \\
\textbf{کاربر ۳ ($U_3$)} & $1$ & $2$ & $5$ & $4$ \\
\textbf{کاربر ۴ ($U_4$)} & $4$ & $3$ & $2$ & $?$ \\
\bottomrule
\end{tabular}
\end{table}

هدف ما پیش‌بینی امتیاز مفقود کاربر ۱ برای \textbf{فیلم ۳} ($\hat{r}_{u_1, i_3}$) است. گام‌های زیر را به صورت دقیق انجام دهید:
\begin{enumerate}
    \item ضریب همبستگی پیرسون (\lr{Pearson Correlation Coefficient}) را بین کاربر ۱ ($U_1$) و سایر کاربران محاسبه نمایید. به صورت خلاصه توضیح دهید چرا در این سناریو، معیار پیرسون نسبت به شباهت کسینوسی (\lr{Cosine Similarity}) ساده، ارزیابی عادلانه‌تری از شباهت کاربران ارائه می‌دهد؟ 
    \item با فرض انتخاب ۲ همسایه نزدیک ($K=2$) بر اساس بالاترین ضریب پیرسون، امتیاز پیش‌بینی‌شده $\hat{r}_{u_1, i_3}$ را با فرمول میانگین وزنی نرمال‌سازی‌شده (\lr{Mean Centering}) محاسبه کنید. (نکته: در صورتی که یکی از همسایگان برتر، خود فاقد امتیاز روی فیلم ۳ باشد، او را کنار گذاشته و همسایه واجد شرایط بعدی را جایگزین کنید).
\end{enumerate}

\item \textbf{ریاضیات فاکتورگیری ماتریس (\lr{Matrix Factorization}) و بهینه‌سازی \lr{SGD}}
\\
تابع هزینه کلي (\lr{Loss Function}) زیر را همراه با مؤلفه‌های بایاس (\lr{Bias}) و تنظیم‌گر (\lr{Regularization}) در نظر بگیرید:
$$L = \sum_{(u,i) \in K} \left( r_{u,i} - (\mu + b_u + b_i + p_u \cdot q_i^T) \right)^2 + \lambda \left( ||p_u||^2 + ||q_i||^2 + b_u^2 + b_i^2 \right)$$
با فرض اینکه خطای پیش‌بینی به صورت $e_{u,i} = r_{u,i} - \hat{r}_{u,i}$ تعریف شده باشد:
\begin{enumerate}
\item \textbf{\textcolor{red}{(امتیازی)}} با توجه به الگوی بهینه‌سازی تصادفی نزول گرادیان (\lr{SGD}) تدریس‌شده در جزوه، روابط نهایی به‌روزرسانی را برای \textbf{بردار فاکتور پنهان آیتم ($q_i$)} و \textbf{پارامتر بایاس کاربر ($b_u$)} با در نظر گرفتن نرخ یادگیری $\eta$ و ضریب تنظیم‌گر $\lambda$ بنویسید. تغییرات هر پارامتر را در هر گام به صورت کیفی تحلیل کنید.    \item مفهوم فاکتورهای پنهان (\lr{Latent Factors}) را بر اساس مثال‌های کلاسی (مانند ابعاد اکشن یا احساسی بودن فیلم‌ها) توضیح دهید  و مشخص کنید چرا این فاکتورها لزوماً با ویژگی‌های صریح کاتالوگ انطباق یک‌به‌یک ندارند؟ 
\end{enumerate}

\item \textbf{ارزیابی معیارهای رتبه‌بندی (\lr{Ranking Metrics})}
\\
یک سیستم توصیه‌گر لیستی با ۵ کاندید برتر از محصولات را به صورت زیر تولید کرده است:
$$Recommended = [Item_A, Item_B, Item_C, Item_D, Item_E]$$
طبق داده‌های تست، تنها اقلام $Item_A$، $Item_C$ و $Item_E$ مرتبط (\lr{Relevant}) هستند و کاربر به آن‌ها امتیازهای واقعی $rel_A = 3$، $rel_C = 5$ و $rel_E = 2$ داده است ($rel_B = 0, rel_D = 0$). با فرض اینکه تعداد کل آیتم‌های مرتبطِ این کاربر در کل بخش تست برابر با \lr{4} باشد:
\begin{enumerate}
    \item مقادیر معیارهای ارزیابی دقت در ۵ بردار اول ($Precision_5$) و فراخوانی در ۵ بردار اول ($Recall_5$) را برای این کاربر به دست آورید.
    \item معیار بهره‌وری تراکمی تخفیف‌خورده نرمال‌سازی‌شده در سطح ۵ ($NDCG_5$) را برای این لیست با محاسبات گام‌به‌گام محاسبه کنید. (برای تشکیل \lr{IDCG}، ابتدا امتیازهای واقعی اقلام مرتبط را به صورت نزولی مرتب کنید) .
\end{enumerate}

\item \textbf{کدنویسی کارگاهی: محاسبه ماتریس شباهت متنی فیلم‌ها}
\\
با استفاده از کتابخانه \texttt{sklearn} در پایتون، یک اسکریپت ساده بنویسید که توضیحات خلاصه داستان سه فیلم زیر را دریافت کند:
\begin{enumerate}
    \item \lr{Inception}: \small{\texttt{"A sci-fi movie about space exploration and time travel."}} 
    \item \lr{Titanic}: \small{\texttt{"A romance drama about two people in love."}}
    \item \lr{Interstellar}: \small{\texttt{"A science fiction story with time travel and space mission."}}
\end{enumerate}
برنامه شما باید با ابزار \texttt{TfidfVectorizer} بردارهای عددی این سه فیلم را استخراج کرده  و در نهایت، ماتریس شباهت کسینوسی متقابل (\lr{Cosine Similarity Matrix}) بین آن‌ها را محاسبه و در خروجی چاپ کند. (از تنظیمات پیش‌فرض برای \texttt{ngram} استفاده کنید).

\end{enumerate}

\vspace{0.5cm}
\noindent
\textbf{تاریخ تحویل:} تا یک روز قبل از آزمون پایان ترم

\end{document}